\chapter{Architecture et Conception du Core}

\section{Architecture hexagonale}
Le module suit la structure hexagonale canonique. Le domaine ne connaît ni les technologies de transport (Web Flux) ni la persistance (bases relationnelles R2DBC). Il communique uniquement à travers des ports d'entrée et de sortie.

\subsection{Couches applicatives et packages}
\begin{center}
\small
\begin{longtable}{@{}p{4cm}p{10cm}@{}}
\toprule
\textbf{Couche / Package} & \textbf{Contenu} \\
\midrule
\code{domain.model} & Agrégat \code{UserAccount}, invariants métier, énumérations de statut. \\
\code{application.port.in} & Cas d'usage entrants (\code{LoginUseCase}, \code{RegisterUserUseCase}, \code{PublicSignUpUseCase}, \code{DiscoverLoginContextsUseCase}). \\
\code{application.port.out} & Ports sortants : \code{UserAccountRepository}, \code{RefreshTokenRepository}, \code{UserOrganizationAccessDirectory}. \\
\code{application.service} & Orchestration : \code{AuthApplicationService}, \code{AuthContextApplicationService}, \code{AuthOidcService}, \code{AuthSharedSessionService}. \\
\code{adapter.in.web} & Contrôleurs WebFlux (\code{AuthController}, \code{AuthOidcController}) et objets de transfert de données (DTO). \\
\code{adapter.out.persistence} & Implémentation de la persistance (R2DBC H2/PostgreSQL). \\
\bottomrule
\end{longtable}
\end{center}

\subsection{Ports sortants et découplage}
Trois ports d'interface principaux découplent \code{auth-core} des autres domaines du Kernel :
\begin{itemize}
  \item \code{UserAccountRepository} : Persistance réactive des comptes.
  \item \code{UserOrganizationAccessDirectory} : Liste des organisations accessibles par un utilisateur.
  \item \code{SignUpContextDirectory} : Résolution des contextes de rattachement lors d'un enregistrement.
\end{itemize}

\section{Modèle de Domaine et Invariants}
L'agrégat central du domaine est la classe immuable \code{UserAccount}. Chaque modification d'état (activation du MFA, changement de plan, vérification d'email) produit une nouvelle instance de l'entité, garantissant la sûreté de la programmation réactive multitâche.

\begin{invbox}[Invariants clés de YowAuth]
\begin{itemize}
  \item L'identifiant (\code{username}) et l'adresse email (\code{email}) sont uniques par Tenant et normalisés en minuscules.
  \item Le hachage du mot de passe (\code{passwordHash}) utilise l'algorithme BCrypt.
  \item L'accès d'un utilisateur dont le MFA est activé ne doit pas renvoyer de session JWT finale avant la validation réussie du défi OTP.
  \item Les comptes suspendus ou désactivés (\code{status = 'DEACTIVATED'|'SUSPENDED'}) voient leur accès systématiquement bloqué lors de l'authentification.
\end{itemize}
\end{invbox}

\section{Diagramme des cas d'utilisation}
Les rôles sont distribués entre trois acteurs :
\begin{itemize}
  \item \textbf{Administrateur IAM} : Enregistre de nouveaux comptes via des accès d'administration sécurisés.
  \item \textbf{Utilisateur final} : Se connecte, configure son MFA, modifie son forfait et consulte son profil.
  \item \textbf{Applications clientes backend} : Réalisent le Token Exchange et valident les jetons de session.
\end{itemize}
